\section{Presentación del proyecto}
\subsection{Nombre y propósito}
La \textbf{Red Descentralizada de Energía Comunitaria} es una aplicación para registrar energía de hogares y gestionar su asignación y pago dentro de una comunidad. Su código está disponible en el repositorio público \href{https://github.com/0x3ef/Decentralized-Community-Energy-Network}{Decentralized-Community-Energy-Network}. El sistema relaciona mediciones, usuarios, créditos SUN y pagos en ETH \citep{software}.

La validación se realizó con sensores simulados y una red local Anvil. Los resultados descritos corresponden a ese entorno de prueba.

\subsection{Problema abordado}
Un hogar puede generar más energía de la que consume mientras otro necesita energía. Para calcular una compensación se requiere conocer cuánto exportó el productor, cuánto consumió el otro hogar y qué cantidad puede asignarse entre ambos.

El proyecto reúne esos datos y aplica reglas para reconocer la energía de cada participante. El proceso comprende la identificación del medidor, la revisión de lecturas, el cálculo por intervalos, la aprobación del lote y la liquidación del pago. La aplicación permite consultar cada etapa y relacionarla con sus registros \citep{readme,settlement}.

\subsection{Objetivo general}
Desarrollar una plataforma que convierta mediciones de energía de hogares en registros de asignación, créditos SUN y pagos entre participantes mediante contratos inteligentes.

\subsection{Objetivos específicos}
\begin{longtable}{L{.27\linewidth}L{.63\linewidth}}
\caption{Objetivos del proyecto y funciones desarrolladas}\label{tab:objetivos}\\
\toprule Objetivo & Función desarrollada \\ \midrule\endfirsthead
\toprule Objetivo & Función desarrollada \\ \midrule\endhead
Registrar participantes & Gestionar usuarios, wallets, hogares, comunidades, nodos y dispositivos.\\
Recibir mediciones & Publicar datos por MQTT y almacenar lecturas autenticadas.\\
Calcular energía & Obtener diferencias entre contadores y agruparlas en intervalos.\\
Asignar excedentes & Distribuir la exportación elegible entre las demandas de la comunidad.\\
Aprobar lotes & Preparar y ejecutar decisiones de gobernanza con una wallet autorizada.\\
Emitir créditos & Entregar SUN por la inyección aprobada.\\
Liquidar pagos & Reservar SUN y ETH, entregar créditos y acreditar ingresos al productor.\\
Registrar compensación & Quemar los SUN correspondientes al consumo aprobado y entregado.\\
Conservar registros & Mantener lecturas, operaciones y estados para su consulta.\\
\bottomrule
\end{longtable}

\subsection{Participantes}
\textbf{Productor.} Aporta exportación elegible, recibe SUN, los deposita en custodia y cobra ETH cuando se liquida una asignación.

\textbf{Consumidor.} Deposita ETH para financiar la energía asignada, recibe SUN y puede utilizarlos para registrar la compensación de su consumo.

\textbf{Operador.} Consulta la operación de la comunidad, los cálculos y las tareas de procesamiento.

\textbf{Oráculo.} Presenta en blockchain los datos del lote calculado por el servidor. Su cuenta tiene permisos para esa tarea.

\textbf{Gobernanza.} Administra las autorizaciones de la comunidad y aprueba los lotes que habilitan la emisión de créditos \citep{blockchain,security}.

\subsection{Fuentes del informe}
El informe se elaboró a partir de la documentación del repositorio, los modelos de datos, los contratos y los resultados de pruebas conservados. Se revisaron las guías de arquitectura, API, IoT, base de datos, liquidación y despliegue. Las cifras de energía y pagos proceden del escenario local de validación. Los resultados de pruebas se presentan por grupo para conservar el contexto de cada ejecución \citep{architecture,database,iot,verification}.
